\section{A tame tower for the prescribed Case III label}
\label{sec:tensor-tower-case-iii}

The Case III argument needs four distinct residual diagonal characters
at one element of the field on which the argument will take place. It
also needs a quadratic field that realizes one particular ratio of
these characters. We construct both fields and the element together.

% BEGIN TFT1
\paragraph{TFT1. Two character values available before changing fields.}
Fix the original algebraic closure $\Omega/\mathbb Q$, prime $\ell$,
coefficient datum $c$, and raw datum $D$ satisfying $\mathrm{Ctx3}$.
Write $O=c.O$, $k=O/\mathfrak m$, $\Delta=G_{\mathbb Q}$ inside
this closure, and use bars for reduction to $k$. The finite character
$\psi$ here is the Galois character of the host Dirichlet character.
Let $m$ be the order of the latter. AX1 and the Case III conditions give
\[
 \ell\ge5,\quad \ell-1\mid 2e,\quad
 m>1\text{ odd},\quad (m,\ell)=1,\quad \psi^m=1.
\]
Indeed $\psi$ is nontrivial by Case III and $m$ divides the odd
part of the host order $o$, which is prime to $\ell$.
Thus no inertia image of $\psi$ has order divisible by $\ell$,
so $D.S_b=\varnothing$. The conductor of $\psi$ divides the level
$N$, prime to $\ell$. The character $\xi$ is the product of an odd
dyadic character of modulus $2^n$ and an auxiliary character of
modulus $p_1$, with $n=n_{\rm III}$, $p_1\ne\ell$ an odd
prime coprime to $N$. Thus $\xi$ and $\psi$ are trivial on
every inertia group above $\ell$. Write $\omega$ for the
Teichm\"uller lift of the residual cyclotomic character.
Reduction is injective on roots of unity of order prime to $\ell$
since $1+\mathfrak m$ is pro-$\ell$.

There are $a\in I_2$ at the specified dyadic place and $u\in\Delta$
with the following values, where $d$ has nontrivial odd order dividing $m$:
\[
 \begin{array}{c|ccc}
      &\xi&\psi&\omega\\ \hline
 a&-1&1&1\\
 u&1&d&1
 \end{array}
\]
At $2$, inertia surjects onto the units modulo $2^n$ because the
dyadic cyclotomic field is totally ramified there. Take an element
acting by $-1$ on $\mu_{2^n}$. It is trivial on the odd conductor
of the auxiliary factor of $\xi$ and on $\mu_\ell$, while the dyadic
factor is odd. Raising to the odd power $m$ gives $a$.
For $u$, pick a unit modulo $N$ on which $\psi$ is nontrivial.
Since $(N,\ell p_1)=1$, the Chinese remainder theorem and the
cyclotomic Galois action give a lift acting by that unit on $\mu_N$
and trivially on $\mu_{\ell p_1}$. Its $2^n$-th power kills the
dyadic factor of $\xi$ and preserves a nontrivial $\psi$-value,
as $m$ is odd. This is $u$.
% END TFT1

% BEGIN TFT2
\paragraph{TFT2. The real field and two distinct quadratic roles.}
Let $S=D.S$, which contains $2$ and $\ell$. Choose a new prime
$q>\ell$ larger than all prime divisors of $2N\ell p_1$ and than
all members of $S$, with $q\equiv-1\pmod{8\prod_{v\in S}v}$.
Dirichlet's theorem applies to this reduced residue class. The
quadratic field $K'=\mathbb Q(\sqrt{-q})\subset\Omega$ has
discriminant $-q$ and associated character $\eta$. Each odd
$v\in S$ splits since $-q\equiv1\pmod v$; $2$ splits since
$-q\equiv1\pmod8$. Inertia at $q$ contains an element $b$
with $\eta(b)=-1$. At this prime $\xi,\psi,\omega$ are all unramified.
The sign character $\eta$ is regarded also as an $O$-valued unit.

Let $\Psi=\chi_\ell^e\psi$ and form the open subgroup
\[
 H=\{g\in\Delta:\ \xi(g)=\eta(g),\quad
                          \bar\Psi(g)=\bar\psi(g)\}.
\]
The two character quotients have finite image and open kernel;
the second is $\bar\omega^e$. Thus $H$ is normal with abelian quotient.
All complex conjugations lie in it: $\xi$ and $\eta$ are odd,
$\Psi$ is even by goodness, and $\psi$ of odd order kills
complex conjugation. For $t$ in any inertia group over $\ell$,
$t^2\in H$. Indeed $\xi(t)=\psi(t)=1$, $\eta(t)^2=1$ and
$\bar\omega(t)^{2e}=1$ since $\ell-1\mid 2e$.

Choose two further, distinct quadratic characters $\lambda,\lambda'$
of negative prime discriminants, by the same progression construction.
For $\lambda$ impose splitting at $S\cup\{q\}$; for $\lambda'$
impose splitting there and at the ramification prime of $\lambda$.
Let $K$ be the field of $\lambda$, put
$J=\ker\lambda\cap\ker\lambda'$, and let $F$ be fixed by
$H\cap\ker(\lambda/\lambda')$. Then
\[
 G_F\subseteq H,\qquad G_F\cap G_K=H\cap J.
\]
Here and below we identify these Galois groups with subgroups
of $\Delta$. The field $F$ is finite abelian and totally real since
both auxiliary characters are odd and $H$ contains complex conjugation.
Its degree is even: $\lambda/\lambda'$ is nontrivial on inertia at the
prime of $\lambda$, and its quadratic field lies in $F$.
The field $K$ is imaginary quadratic and splits every prime of $S$.
Put $L=FK$. At the prime of $\lambda$, inertia inside $G_F$
has $\lambda=\lambda'=1$; away from that prime $\lambda$ is unramified.
Consequently $L/F$ is unramified at every finite place.
Every square in $\Delta$ lies in $J$, and $J$ contains every absolute
decomposition group over $S\cup\{q\}$. In particular
$t^2\in G_L$ for all inertia elements $t$ over $\ell$.
% END TFT2

% BEGIN TFT3
\paragraph{TFT3. One separating element for either allowed label.}
Fix the supplied Boolean label $\mathrm{pr}$ satisfying
$\mathrm{LabOk}(D,\mathrm{pr})$. Set $\epsilon=1$ for the false
label and $-1$ for the true label. Both $a$ and $b$ of TFT1--2
belong to $J$ by the splitting at $2,q$, and $u^2\in J$.
Put $g_0=a b u^2$. Then
\[
 \xi(g_0)=\eta(g_0)=-1,\qquad
 \omega(g_0)=1,\qquad \psi(g_0)=d^2 .
\]
In fact $\eta(a)=1$ by splitting at $2$, and $\xi,\omega,\psi$
are trivial on $b$. Thus $g_0\in H\cap J=G_L$. Write
$z=\bar d^{\,2\epsilon}$. This is of nontrivial odd order, since
reduction preserves the order of $d$ and its order is odd.
In particular $z\ne 1,-1$.

For the relabelled four diagonal characters $\chi_i$, their quotients
$\chi_0/\chi_i$ in order are $1,\alpha,\Psi^\epsilon,
\alpha\Psi^\epsilon$, where $\alpha=\chi_\ell\omega^{-1}\xi$.
This follows directly from the two permutations of the diagonal.
On $G_F$, $\bar\alpha=\bar\xi$ and $\bar\Psi=\bar\psi$.
Consequently the residual root at $w=(h,k')$ is
\[
 r_w=\bar\xi^h(\bar\psi^\epsilon)^{k'},\qquad
 r_w(g_0)=(-1)^h z^{k'}.
\]
At $g_0$ the reduced four quotients are $1,-1,z,-z$, all distinct
as $\ell\ge5$. Thus the four diagonal values themselves are distinct.
The eight roots are $\pm(1,0),\pm(0,1),\pm(1,1),\pm(-1,1)$.
Their values are among $-1,z^{\pm1},-z^{\pm1}$ and are all
nontrivial. The root $(1,1)$ has value $-z$ of order
$2\operatorname{ord}(z)\ge6$, so no power from one through four kills it.

Suppose a quadratic character $\nu$ of $G_L$ is trivial on
$G_L\cap\ker\omega$ and multiplication by $\nu$ permutes the four
restricted diagonal characters in the sense of the source. Apply this
condition to $\bar\chi_0$. The resulting equality
$\nu\bar\chi_0=\bar\chi_j$ at $g_0$ has $\nu(g_0)=1$, so
separation gives $j=0$. The same equality on all of $G_L$ gives
$\nu=1$. All three character tests use the same field and the
same $g_0$ for the initially prescribed label.
% END TFT3

% BEGIN TFT4
\paragraph{TFT4. Inertia and the local alternatives on this tower.}
For every place of $\Omega$ over $\ell$, the mod-$\ell$ cyclotomic
image on rational inertia is $\mathbb F_\ell^\times$, since
$\mathbb Q_\ell(\mu_\ell)/\mathbb Q_\ell$ is totally ramified of
degree $\ell-1$. Choose an inertia element $t$ mapping to a generator.
Its square has $\bar\omega(t^2)\ne1$ because $\ell-1\ge4$,
and $t^2\in G_L$ by TFT2. The relative inertia is the intersection
of absolute inertia with $G_F$, and $\xi(t)=\psi(t)=1$. Hence
$r_w(t^2)=1\ne\bar\omega(t^2)$ for \emph{every} $w\in\mathbb Z^2$.
This gives both the cyclotomic exclusion on $G_L$ and, at every
place over $\ell$, the first alternative of $\mathrm{locL}$ on
relative inertia. At those places all inertia elements, on squaring,
fix $F$; the uniform exponent $2$ is prime to $\ell$ and gives
exactly $\mathrm{TameL}(\ell,F)$.

The original diagonal residual representation $\bar r_0$ is trivial
at an element precisely when $\bar\alpha=\bar\Psi=1$: its diagonal
characters are $\alpha\Psi,\Psi,\alpha,1$ before relabelling.
On relative inertia over $\ell$, $\bar\alpha=\bar\xi=1$ and
$\bar\Psi=\bar\psi=1$ by TFT2. Away from $\ell$,
$\omega$ is unramified, hence $\bar r_0=1$ on an inertia element
implies $1=\bar\Psi=\bar\psi$ there. Since $\psi^m=1$ and
$(m,\ell)=1$, this gives $\psi=1$ in $O$. Finally, $\psi^m=1$
holds on all of $G_F$, in particular on its full cyclotomic kernel.
These are the integral and residual inertia clauses of $\mathrm{Ok}$.
% END TFT4

% BEGIN TFT5
\paragraph{TFT5. The original tensor-factor tower exists.}
Use $F,K,L$ from TFT2 and $g_0$ from TFT3. The compositum $L=FK$
is CM, quadratic over its totally real subfield $F$, since $K$
is imaginary quadratic. For every prime of $F$ above $\ell$ take
one prolongation to $\Omega$ whose restriction to $K$ is that of
the prescribed valuation $D.\mathrm{pl}(\ell)$. Such a choice is
possible because the two primes of $K$ over $\ell$ split after
base change to each completion $F_v$, and the involution of $L/F$
interchanges the two local factors. Absolute Galois over $F$ acts
transitively on the prolongations of each $v$. These finitely many
representatives therefore give exactly the original orbit and
place-on-$K$ requirements. Together with the supplied label these
data define the tower $T$ inside the original $\Omega$.

For $T.\mathrm{Ok}$, the label hypothesis is the one supplied, and
TFT2 gives even degree, unramified $L/F$ and splitting of $S$ in
$K$. TFT3 gives separation at $T.g_0$, nontriviality of each
root, the order test and exclusion of the quadratic self-twists.
TFT4 gives the remaining root, inertia and local tests.
For $\mathrm{Four}(T)$, TFT1 gives $\ell\ge5$, the empty bad set
and a prime-to-$\ell$ exponent for the original $\psi$.
Its quadratic witness is $K'=\mathbb Q(\sqrt{-q})$ of TFT2.
It is totally complex of degree two and $\ell\nmid\operatorname{disc}(K')$.
The decomposition groups above $\ell$ and all primes of $D.S$
fix $K'$. For $g\in G_F$ the root $w_\alpha=(1,0)$ has value
$\bar\xi(g)=\bar\eta(g)$, hence equals $1$ on $G_F\cap G_{K'}$
and $-1$ on its complement. This supplies the two labelled
restrictions of the original $\mathrm{QuadOn}$ record.
The same tower is tame by TFT4. We have proved, for every original
$D$ in Case III and every permitted $\mathrm{pr}$, the existence of
\[
 T\quad\text{with}\quad T.\mathrm{pr}=\mathrm{pr},\quad
 T.\mathrm{Ok},\quad \mathrm{Four}(T),\quad
 \mathrm{TameL}(\ell,T.F).
\]
This supplies the initial tower used in CJP3. Producing the seeded
pair on a later tower belongs to the separate rank-four reduction.
% END TFT5
